\documentclass[12pt, a4paper]{article}
\usepackage[utf8]{inputenc}
\usepackage{amsmath, amssymb, amsthm, amsfonts,amsthm,tikz-cd}
\usepackage{mathrsfs}
\usepackage{CJKutf8}
\usepackage{bm}
\usepackage[margin=1in]{geometry}
\theoremstyle{remark}
\newtheorem*{remark}{Remark}
\newenvironment{lemma}[1]{
    \noindent \textbf{Lemma #1.} \itshape
}{
    \vspace{1em}
}

% 重新定义 solution 环境，保留标识并在末尾留出 8cm 空白
\newenvironment{solution}
  {\paragraph{\textit{Solution:}}\par\medskip}
  {\hfill$\blacksquare$\vspace{0.1cm}}

\title{Abstract Algebra III: Assignment - Week 4}
\author{
    \begin{CJK*}{UTF8}{gbsn}
    钟皓宇 (12533556)
    \end{CJK*}
}
\date{\today}

\begin{document}

\maketitle
\section*{Problem 1}
Let $M, N, M_i (i \in I)$ and $N_j (j \in J)$ be $R$-modules. Prove the following isomorphisms of abelian groups.
\begin{enumerate}
    \item[(1)] $\operatorname{Hom}_R\left(\bigoplus_{i\in I} M_i, N\right) \simeq \prod_{i\in I} \operatorname{Hom}_R(M_i, N)$;
    \item[(2)] $\operatorname{Hom}_R\left(M, \prod_{j\in J} N_j\right) \simeq \prod_{j\in J} \operatorname{Hom}_R(M, N_j)$.
\end{enumerate}
\begin{solution}
\textbf{(1)}
Let $X = \operatorname{Hom}_R\left(\bigoplus_{i\in I} M_i, N\right)$ and $Y_i = \operatorname{Hom}_R(M_i, N)$. 
Let $G$ be an arbitrary abelian group, and suppose we are given a family of group homomorphisms $\varphi_i: G \to Y_i$ for each $i \in I$. 
For fixed $g\in G$, the family $\{\varphi_i(g)\}_{i \in I}$ is a collection of $R$-module homomorphisms from $M_i$ to $N$.
Then by the universal property of the direct sum $\bigoplus_{i \in I} M_i$ in $R$-Mod, there exists a unique $R$-module homomorphism $f_g$
such that 
\begin{align*}
    f_g \circ \iota_i = \varphi_i(g) \quad \text{for all } i \in I.
\end{align*}
where $\iota_i$ are the canonical injections. Define $\Phi: G \to X$ by setting $\Phi(g) = f_g$. I claim that $\Phi$ is the unique group 
homomorphism such that the diagram
\[
\begin{tikzcd}[ampersand replacement=\&, row sep=large, column sep=large]
    G \arrow[d, "\varphi_i"'] \arrow[dr, "\Phi", dashed] \& \\
    Y_i \& X \arrow[l, "\pi_i"]
\end{tikzcd}
\]
commutes where $\pi_i=\text{Hom}_R(\iota_i,N)$ is the canonical projection induced by contravariant functor $\text{Hom}_R(-,N)$. Let $g, h \in G$. Since $\varphi_i(g+h)= \varphi_i(g)+\varphi_i(h)$ we obtain that
\begin{align*}
    f_{g+h}\circ \iota_i = \varphi_i(g)+\varphi_i(h) = (f_{g}+f_{h})\circ \iota_i.
\end{align*}
Then $f_{g+h} = f_g+f_h$, or equivalently $\Phi(g+h)=\Phi(g)+\Phi(h)$, by the uniqueness condition of the direct sum's universal property.  Therefore, $\Phi$ is a 
group homomorphism. To prove the uniqueness of $\Phi$, we suppose
there is another group homomorphism $\Psi: G \to X$ satisfying $(\Psi(g)) \circ \iota_i = \varphi_i(g)$ for all $i \in I$ and all $g \in G$. For any fixed $g$, both $R$-module homomorphisms $\Psi(g)$ and $\Phi(g)$ pull back along $\iota_i$ to the same maps $\varphi_i(g)$. 
By the universal property of the direct sum, we must have $\Psi(g) = \Phi(g)$ for all $g$, so $\Psi = \Phi$.
The uniqueness and existence of $\Phi$ yields $\operatorname{Hom}_R\left(\bigoplus_{i\in I} M_i, N\right) \simeq \prod_{i\in I} \operatorname{Hom}_R(M_i, N)$.

\textbf{(2)} Let $X=\operatorname{Hom}_R\left(M, \prod_{j\in J} N_j\right)$ and $Y_j=\operatorname{Hom}_R(M, N_j)$. Let $G$ be an arbitrary abelian group, and suppose we are given a family of group homomorphisms
$\varphi_j: G \to Y_j$ for each $j \in J$. Then for fixed $g\in G$ the family of $R$-module homomorphisms $\{\varphi_j(g):M\to N_j\}_{j\in J}$ uniquely induced an $R$-module homomorphism
$f_g:M \to \prod_{j\in J} N_j$ such that $\varphi_j(g)= p_j\circ f_g$. Define the map $\Phi:G \to \operatorname{Hom}_R\left(M, \prod_{j\in J} N_j\right)$ by setting $\Phi(g) = f_g$. I claim $\Phi$ is the unique group
homomorphism such that the diagram 
\[
\begin{tikzcd}[ampersand replacement=\&, row sep=large, column sep=large]
    G \arrow[d, "\varphi_j"'] \arrow[dr, "\Phi", dashed] \& \\
    Y_j \& X \arrow[l, "\pi_j"]
\end{tikzcd}
\]
commutes where $\pi_j=\text{Hom}_R(M,p_j)$ is the canonical projection induced by covariant functor $\text{Hom}_R(M,-)$. Let $g, h \in G$. Since $\varphi_j(g+h)= \varphi_j(g)+\varphi_j(h)$ we obtain that
\begin{align*}
    p_j\circ f_{g+h} = \varphi_j(g)+\varphi_j(h) = p_j\circ (f_{g}+f_{h}).
\end{align*}
It implies $\Phi(g+h)=\Phi(g)+\Phi(h)$ because of the uniqueness condition of the direct product's universal property.
Therefore, $\Phi$ is a group homomorphism. To prove the uniqueness of $\Phi$, we suppose
there is another group homomorphism $\Psi: G \to X$ satisfying $p_j\circ (\Psi(g))  = \varphi_j(g)$ for all $j \in J$ and all $g \in G$. For any fixed $g$, both $R$-module homomorphisms $\Psi(g)$ and $\Phi(g)$ push forward along $p_j$ to the same maps $\varphi_j(g)$. 
By the universal property of the direct product, we must have $\Psi(g) = \Phi(g)$ for all $g$, so $\Psi = \Phi$.
The uniqueness and existence of $\Phi$ yields $\operatorname{Hom}_R\left(M, \prod_{j\in J} N_j\right) \simeq \prod_{j\in J} \operatorname{Hom}_R(M, N_j)$.
\end{solution}
\begin{remark}
    The isomorphism established in \textbf{(1)} is the direct consequence of the \textbf{Hom-Hom adjunction} and the fact that right adjoints preserve limits. Specifically, 
    the contravariant functor $\operatorname{Hom}_R(-, N)$ is the right adjoint of contravariant functor $\operatorname{Hom}_{\textbf{Ab}}(-, N)$. 
    More explicitly, there is a natural isomorphism 
    \begin{align*}
        \text{Hom}_R(M,\text{Hom}_\textbf{Ab}(A,N)) \cong \text{Hom}_\textbf{Ab}(A,\text{Hom}_R(M,N)),
    \end{align*}
    which holds for both left and right $R$-module categories.
    Similarily, \textbf{(2)} is the direct consequence of the \textbf{Tensor-Hom adjunction}. Specifically, the covariant functor
    $\text{Hom}_R(M,-)$ is the right adjoint of the covariant functor $-\otimes_\mathbb{Z} M$ if we consider right $R$-module category and $M \otimes_\mathbb{Z} - $ for the left.
\end{remark}


\section*{Problem 2}
(\textit{Five Lemma}) Suppose that the following diagram of homomorphisms of $R$-modules
\[\begin{tikzcd}
    M_1 \arrow[r,"f_1"] \arrow[d, "\varphi_1"] & M_2 \arrow[r,"f_2"] \arrow[d, "\varphi_2"] & M_3 \arrow[r,"f_3"] \arrow[d, "\varphi_3"] & M_4 \arrow[r,"f_4"] \arrow[d, "\varphi_4"] & M_5 \arrow[d, "\varphi_5"] \\
    N_1 \arrow[r,"g_1"] & N_2 \arrow[r,"g_2"] & N_3 \arrow[r,"g_3"] & N_4 \arrow[r,"g_4"] & N_5
\end{tikzcd}\]
is commutative and has exact rows. Prove that
\begin{enumerate}
    \item[(1)] if $\varphi_1$ is surjective, and $\varphi_2$ and $\varphi_4$ are injective, then $\varphi_3$ is injective;
    \item[(2)] if $\varphi_5$ is injective, and $\varphi_2$ and $\varphi_4$ are surjective, then $\varphi_3$ is surjective.
\end{enumerate}

\begin{solution}
    \textbf{(1)} Suppose $a_3\in M_3$ such that $\varphi_3(a_3) = 0$, then $\varphi_4\circ f_3(a_3) = g_3\circ \varphi_3(a_3)=0$ implies $f_3(a_3)=0$
    since $\varphi_4$ is injective. Therefore, $\text{Im}(f_2)=\ker(f_3)$ implies that there exists an element $a_2\in M_2$ such that $f_2(a_2)=a_3$.
    Let $b_2 = \varphi_2(a_2)$, then $g_2(b_2) = \varphi_3\circ f_2(a_2)= \varphi_3(a_3)=0$ implies there exists $b_1\in N_1$ such that $g_1(b_1)=b_2$. 
    Because $\varphi_1$ is surjective then there exists $a_1$ such that $\varphi_1(a_1)=b_1$. I claim that $a_2=f_1(a_1)$ and therefore $a_2\in \text{Im}(f_1)=\ker(f_2)$ implies
    $a_3=f_2(a_2)=0$ and thus $\varphi_3$ is injective. To prove it, we can consider $\varphi_2(a_2-f_1(a_1))=\varphi_2(a_2) - g_1\circ \varphi_{1}(a_1) = b_2-b_2=0$  which 
    implies $a_2-f_1(a_1)=0$ because of injectivity of $\varphi_2$.
    
    \textbf{(2)} Suppose $b_3$ be arbitrary elements in $N_3$ respectively. Let $b_4=g_3(b_3)$. Because $\varphi_4$ is surjective 
    then we can let $a_4 \in M_4$ such that $\varphi_4(a_4)=b_4$. Then $\varphi_5\circ f_4(a_4)=g_4\circ \varphi_4(a_4) = g_4\circ g_3(b_3) = 0$ implies 
    $f_4(a_4)=0$ because of injectivity of $\varphi_5$. Then we can suppose $a_3\in M_3$ such that $f_3(a_3)= a_4$. Then $g_3(b_3-\varphi_3(a_3))= b_4-\varphi_4(a_4)=0$ implies 
    there exists $b_2\in N_2$ such that $g_2(b_2) = b_3-\varphi_3(a_3)$. Then the surjectivity of $\varphi_2$ implies there exists $a_2$ such that $\varphi_2(a_2)=b_2$.
    Since diagram commutes, we obtain that $\varphi_3\circ f_2(a_2) = b_3-\varphi_3(a_3)$, or equivalently, $b_3 = \varphi_3(f_2(a_2)+a_3)\in \text{Im}(\varphi_3)$ which implies $\varphi_3$ is 
    surjective.
\end{solution}

\section*{Problem 3}
(\textit{Snake Lemma}) Suppose that the following diagram of homomorphisms of $R$-modules
\[\begin{tikzcd}
    0 \arrow[r] & M' \arrow[r, "\alpha"] \arrow[d, "\varphi'"] & M \arrow[r, "\beta"] \arrow[d, "\varphi"] & M'' \arrow[r] \arrow[d, "\varphi''"] & 0 \\
    0 \arrow[r] & N' \arrow[r, "\mu"] & N \arrow[r, "\nu"] & N'' \arrow[r] & 0
\end{tikzcd}\]
is commutative and has exact rows. Prove that
\[ 0 \to \ker\varphi' \xrightarrow{\alpha'} \ker\varphi \xrightarrow{\beta'} \ker\varphi'' \xrightarrow{\delta} \operatorname{coker}\varphi' \xrightarrow{\mu'} \operatorname{coker}\varphi \xrightarrow{\nu'} \operatorname{coker}\varphi'' \to 0 \]
is a long exact sequence, where $\alpha', \beta', \mu', \nu'$ are homomorphisms induced by $\alpha, \beta, \mu, \nu$, respectively. (Hint: for any $x \in \ker\varphi''$, take $y \in M$ such that $\beta(y) = x$, and take $z \in N'$ such that $\mu(z) = \varphi(y)$. Define $\delta(x) = \overline{z} = z + \operatorname{im}\varphi'$)
\begin{solution}
    \textbf{$\bm{\alpha'}$ is well-defined and injective}. Suppose $m'\in\ker(\varphi')$, since diagram commutes, we obtain that $\varphi\circ \alpha(m')=\mu\circ\varphi'(m')=0$ which implies $\alpha(m')\in\ker(\varphi)$. 
    Therefore, $\alpha'= \alpha|_{\ker(\varphi')}:\ker(\varphi')\to \ker(\varphi)$ is well-defined. Hence, $\alpha'$ is injective since $\alpha$ is injective.
    \textbf{$\bm{\beta'}$ is well-defined and $\bm{\textbf{Im}(\alpha')=\ker(\beta')}$.}  Let $m\in\ker(\varphi)$, then $\varphi''\circ\beta(m)=\nu\circ \varphi(m)=0$ yields well-definedness. Suppose $m'\in\ker(\varphi')$ then $\beta'\circ\alpha'(m')=\beta\circ\alpha(m')=0$ where $\beta\circ\alpha=0$ because of 
    the exactness of the first row. Therefore, we obtain that $\text{Im}(\alpha')\subseteq\ker(\beta')$. Conversely, if $m\in\ker(\beta')$, then $\beta(m)=\beta'(m)=0$. By exactness, there exists $m'$ such that $\alpha(m')=m$. Hence, $m\in\ker(\varphi)$ implies that $\mu\circ\varphi'(m')= \varphi\circ\alpha(m')= \varphi(m)=0$. 
    Because $\mu$ is injective, then we obtain that $m'\in\ker(\varphi')$. Therefore, for an arbitrary element $m\in\ker(\beta')$ there exists $m'\in\ker(\varphi')$ such that $m=\alpha'(m')$ meaning that $\ker(\beta')\subseteq\text{Im}(\alpha')$.
    \textbf{$\bm{\delta}$ is well-defined.}
    Because $\beta$ is surjective then for any $x\in \ker(\varphi'')$ there exists an element $y\in M$ such that $\beta(y)=x$.
    Since $x\in\ker(\varphi'')$, we obtain that $\nu\circ\varphi(y)= \varphi''(x)=0$, which implies $\varphi(y)\in\ker(\nu)$. By exactness, 
    there exists an element $z\in N'$ such that $\mu(z)=\varphi(y)$.
    Let $y_0,y_1\in M $ be arbitrary two elements such that $\beta(y_0)=\beta(y_1)=x$. Let $z_0,z_1\in N'$ be arbitrary two elements such that $\mu(z_0)=\varphi(y_0)$ and $\mu(z_1)=\varphi(y_1)$.
    Because $\beta(y_1-y_0)= 0$, then there exists an element $m'\in M'$ such that $\alpha(m')=y_1-y_0$. Therefore, 
    \begin{align*}
        \mu(\varphi'(m')-z_1+z_0)&= \mu\circ \varphi'(m')-\mu(z_1-z_0) \\
        &= \varphi\circ\alpha(m')-\mu(z_1-z_0) \\
        &= \varphi(y_1-y_0) - \mu(z_1-z_0)=0
    \end{align*}
    implies $\varphi'(m')=z_1-z_0$ since $\mu$ is injective. Equivalently, we obtain that $\overline{z_1-z_0}=0$ which yields the well-definedness of $\delta$.
    \textbf{$\bm{\textbf{Im}(\beta')=\ker(\delta)}$}. Suppose $m\in\ker(\varphi)$, then we can select $z=0\in N'$ because $\mu(z)=\varphi(m)=0$ which implies $\delta\circ\beta'(m)=0$.
    Therefore, $\text{Im}(\beta')\subseteq \ker(\delta)$. Suppose $m''\in\ker(\delta)$, then there exists $m\in M$ and $n'\in N'$ such that $\beta(m)=m''$ and $\varphi(m)=\mu(n')$.
    Hence $\delta(m'')=0$ implies there exists an element $m'\in M'$ such that $\varphi'(m')=n'$.  I claim that $m-\alpha(m')\in \ker(\varphi)$ and thus $\beta'(m-\alpha(m'))=\beta(m)-\beta\circ\alpha(m')=m''$ implies 
    that $\ker(\delta)\subseteq\text{Im}(\beta')$. To see that, we consider $\varphi(m-\alpha(m')) = \varphi(m) -\mu\circ\varphi'(m') = \varphi(m)-\mu(n')=0$.
    \textbf{$\bm{\mu'}$ is well-defined and $\bm{\textbf{Im}(\delta)=\ker(\mu')}$.} Since the diagram commutes, $\mu\circ\varphi' = \varphi\circ\alpha$ implies $\mu(\text{Im}(\varphi'))\subseteq \text{Im}(\varphi)$ which 
    yields well-definedness. Let $m''$ be an arbitrary element in $\ker(\varphi'')$ then there exists $m\in M$ and $n'\in N'$ such that $\beta(m) = m''$ and $\varphi(m)=\mu(n')$. Then we have $\mu'\circ\delta(m'') = \mu'(n'+\text{Im}(\varphi')) = \mu(n') + \text{Im}(\varphi)=0\in\text{coker}(\varphi)$,
    which implies $\text{Im}(\delta)\subseteq \ker(\mu')$. For an arbitrary element $\overline{n'}\in\ker(\mu')$, then we have 
    $\mu(n')\in\text{Im}(\varphi)$, which implies there exists an element $m\in M$ such that $\varphi(m)=\mu(n')$. Let $m'' = \beta(m)$ and I claim that $m''\in \ker(\varphi'')$
    and thus $\delta(m'') = \overline{n'}$ implies $\ker(\mu')\subseteq \text{Im}(\delta)$. To see that, just consider $\varphi''(m'')=\nu\circ\varphi(m) = \nu\circ\mu(n')=0$.
    \textbf{$\bm{\nu'}$ is well-defined and $\bm{\ker(\nu')=\textbf{Im}(\mu')}$.} Similarily, because diagram commutes, then $\nu\circ\varphi = \varphi''\circ\beta$ implies that $\nu(\text{Im}(\varphi))\subseteq \text{Im}(\varphi'')$
    meaning that $\nu'$ is well-defined. Let $\overline{n'}$ be an arbitrary element in $\text{coker}(\varphi')$, then $\nu'\circ\mu'(\overline{n'}) = \nu'(\mu(n')+\text{Im}(\varphi))= \nu\circ\mu (n') + \text{Im}(\varphi'') = 0\in \text{coker}(\varphi'')$.
    Therefore, $\text{Im}(\mu')\subseteq \ker(\nu')$. Let $\overline{n}$ be an arbitrary element such that $\nu'(\overline{n})=0$, equivalently $\nu(n)\in \text{Im}(\varphi'')$. Then there exists
    an element $m''\in M''$ such that $\varphi''(m'')=\nu(n)$. Since $\beta$ is surjective, then there exists an element $m\in M$ such that $\beta(m)=m''$. I claim that 
    $n-\varphi(m)\in\ker(\nu)$ and thus there exists an element $n'\in N'$ such that $\mu(n')= n-\varphi(m)$, which implies that $\mu'(\overline{n'}) = \overline{n}$ meaning that $\ker(\nu')\subseteq\text{Im}(\mu')$.
    To see that, we consider $\nu(n-\varphi(m)) = \nu(n)-\varphi''\circ\beta(m)=\nu(n)-\varphi''(m)=0$.
    \textbf{$\bm\nu'$ is surjective.} It is directly because $\nu'$ is well-defined and induced by surjection $\nu$. Explicitly, for any arbitrary element $\overline{n''}\in\text{coker}(\varphi'')$, there exists an element $n\in N$ such that 
    $\nu(n)=n''$. Since $\nu(\text{Im}(\varphi))\subseteq \text{Im}(\varphi'')$ we have discussed, then we have $\nu'(\overline{n})=n''+\text{Im}(\varphi'')=\overline{n''}$. 
\end{solution}

\pagebreak



\section*{Lemmas for Problem 3}

\begin{lemma}{A}
Let $f: M \to M''$ and $g: K \to M''$ be $R$-module homomorphisms. Let $P = M \times_{M''} K$ be the pullback of $f$ along $g$, defined explicitly by the subset of the direct sum:
\[ P = \{ (m, k) \in M \times K \mid f(m) = g(k) \} \]
equipped with the canonical projections $\pi_M: P \to M, (m, k) \mapsto m$ and $\pi_K: P \to K, (m, k) \mapsto k$. Then the following properties hold:
\begin{enumerate}
    \item If $f$ is surjective, then the projection $\pi_K$ is surjective.
    \item If $f$ is injective, then the projection $\pi_K$ is injective.
\end{enumerate}
\end{lemma}

\begin{proof}
\textbf{(1)} Suppose $f$ is surjective. Let $k \in K$ be an arbitrary element.
Because $f$ is surjective, then there exists an element $m\in M$ such that $f(m)=g(k)$, which implies $(m,k)\in P$ then $k=\pi_{K}(m,k)$ implies surjectivity.


\textbf{(2)} Suppose $f$ is injective. Suppose $(m,k)\in P$ such that $\pi_K(m,k)=0$. 
 Then we obtain that $f\circ \pi_M(m,k) = f(m) =g(k)= g\circ \pi_K(m,k) = 0$. The injectivity of $f$ implies $m=0$. Hence, $k=\pi_{K}(m,k)=0$ implies the injectivity of $\pi_K$.
\end{proof}

\begin{lemma}{B}
Let $\mu: N' \to N$ and $h: N' \to L$ be $R$-module homomorphisms. Let $Q = N \oplus_{N'} L$ be the pushout of $\mu$ along $h$, defined explicitly as the quotient module of the direct sum:
\[ Q = (N \oplus L) / W \]
where $W = \{(\mu(n'), -h(n')) \mid n' \in N'\}$. The induced maps are given by $\bar{h}: N \to Q, n \mapsto \overline{(n, 0)}$ and $i_L: L \to Q, l \mapsto \overline{(0, l)}$. Then the following properties hold:
\begin{enumerate}
    \item If $\mu$ is surjective, then the induced map $i_L$ is surjective.
    \item If $\mu$ is injective, then the induced map $i_L$ is injective.
\end{enumerate}
\end{lemma}
\begin{proof}
\textbf{(1)} Suppose $\mu$ is surjective. It suffices to prove that any pair with form $\overline{(n,0)}$ has preimage because $\overline{(n,l)} = \overline{(n,0)}+\overline{(0,l)}$ and $\overline{(0,l)}$ has preimage.
Because $\mu$ is surjective, then there exists $n'\in N'$ such that $\mu(n')=n$. I claim that $\overline{(n,0)} = \overline{(0,h(n'))}= i_L(h(n'))$ and thus surjectivity of $i_L$ holds.
To see that, we consider $\overline{(n,0)}-\overline{(0,h(n'))} = \overline{(n,-h(n'))}=0$ since $n=\mu(n')$.

\textbf{(2)} Suppose $\mu$ is injective. Let $l\in L$ be an arbitrary element such that $i_L(l)= \overline{(0,l)}=0$, or equivalently there exists $n'\in N'$ such that $\mu(n')=0$ and $-h(n')=l$.
The injectivity of $\mu$ implies $n'=0$ and hence $l=-h(n')=0$. Therefore, $i_L$ is injective.
\end{proof}
\end{document}